% SEGMENT OLP-0625-S01
% Part: history
% Chapter: biographies
% Section: kurt-goedel

\documentclass[../../../include/open-logic-section]{subfiles}

\begin{document}

\olfileid{his}{bio}{god}

\olsection{கர்ட் கோடல்}

\olphoto{goedel-kurt}{கர்ட் கோடல்}

கர்ட் கோடல் (\textsc{ger}-dle) 1906 ஏப்ரல் 28 அன்று ஆஸ்திரிய--ஹங்கேரியப்
பேரரசின் பிருன் நகரில் பிறந்தார்; அது இப்போது செக் குடியரசின் பிர்னோ ஆகும்.
சிறு வயதிலேயே ஆர்வமுள்ள, புத்திசாலித்தனமான குழந்தையாக இருந்ததால்,
குடும்பத்தினர் அவரை ``Der kleine Herr Warum'' (``ஏன்'' என்று கேட்கும் சிறிய
ஐயா) என்று அடிக்கடி அழைத்தனர். தொடக்கப் பள்ளியிலிருந்தே பாடங்களில் சிறந்து
விளங்கினார்; கணிதத்தில் மட்டுமே அவருக்கு மிக உயர்ந்த மதிப்பெண்ணைவிடக்
குறைந்த மதிப்பெண் கிடைத்தது. உடல்நலம் குன்றியதால் பள்ளிக்கு அடிக்கடி
வரமுடியாமல் போனது; உடற்கல்வியிலிருந்தும் விலக்கு பெற்றார். சிறுவயதில்
அவருக்கு வாதக் காய்ச்சல் ஏற்பட்டதாகக் கண்டறியப்பட்டது. அது தமது இதயத்தை
நிரந்தரமாகப் பாதித்ததாக அவர் வாழ்நாள் முழுவதும் நம்பினார்; மருத்துவ
மதிப்பீடு அதற்கு மாறாக இருந்தது.

% SEGMENT OLP-0625-S02
கோடல் 1924-இல் வியன்னா பல்கலைக்கழகத்தில் படிக்கத் தொடங்கி, 1929-இல்
முனைவர் படிப்பை முடித்தார். முதலில் இயற்பியல் படிக்க எண்ணியிருந்தார்;
ஆனால் மெய்யியலாளர் ரூடால்ஃப் கார்னாப்பின் தாக்கம் உள்ளிட்ட காரணங்களால்
விரைவில் கணிதம், குறிப்பாகத் தருக்கவியல், மீது ஆர்வம் கொண்டார். ஹான்ஸ்
ஹானின் மேற்பார்வையில் எழுதிய ஆய்வுரையில், சமத்துவம் அடங்கிய முதல்தரப்
பயனிலைத் தருக்கத்தின் நிறைவுத் தேற்றத்தை நிறுவினார் \citep{Godel1929}.
ஓராண்டுக்குப் பின்னரே தமது மிகப் புகழ்பெற்ற முடிவுகளான முதல், இரண்டாம்
முழுமையின்மைத் தேற்றங்களைப் பெற்றார் (\citealt{Godel1931} இல்
வெளியிடப்பட்டவை). வியன்னாவில் இருந்தபோது, அறிவியல் நோக்குடைய
மெய்யியலாளர்களின் வியன்னா வட்டத்தில் தீவிரமாகப் பங்கேற்றார். அதில்
இடம்பெற்ற கார்னாப்பின் பணி, கோடலின் முடிவுகளால் குறிப்பாகத் தாக்கம்
பெற்றது.

% SEGMENT OLP-0625-S03
1938-இல் கோடல் அடேல் நிம்புர்ஸ்கியை மணந்தார். அவர் கோடலைவிட ஆறு
வயது மூத்தவர், ஏற்கெனவே மணமுறிவு பெற்றவர், இரவு விடுதியில் நடனக்
கலைஞராகப் பணியாற்றியவர் என்பதால் கோடலின் பெற்றோர் மகிழ்ச்சியடையவில்லை.
ஆனால் சமூக அழுத்தங்கள் கோடலைப் பாதிக்கவில்லை; அவர் இறக்கும்வரை இருவரும்
மகிழ்ச்சியாக மணவாழ்க்கை நடத்தினர்.

% SEGMENT OLP-0625-S04
1938-இல் நாசி ஜெர்மனி ஆஸ்திரியாவை இணைத்துக்கொண்ட பின்னர், கோடலும்
அடேலும் அமெரிக்க ஐக்கிய நாடுகளுக்குக் குடிபெயர்ந்தனர். நியூ ஜெர்சியின்
பிரின்ஸ்டனில் உள்ள மேம்பட்ட ஆய்வுகள் நிறுவனத்தில் அவர் பணியேற்றார்.
உள்முகத் தன்மையும் வழக்கத்திற்கு மாறான இயல்பும் கொண்டிருந்தபோதிலும்,
பிரின்ஸ்டனில் அவர் கழித்த காலம் கூட்டுப் பணிகளும் பயனுள்ள ஆய்வுகளும்
நிறைந்ததாக இருந்தது. கணக்கோட்பாடு, மெய்யியல், இயற்பியல் ஆகிய துறைகளில்
கட்டுரைகளை வெளியிட்டார். குறிப்பாக, அதே நிறுவனத்தில் பணியாற்றிய
ஆல்பர்ட் ஐன்ஸ்டைனுடன் நெருங்கிய நட்பு கொண்டார்.

% SEGMENT OLP-0625-S05
பிற்காலத்தில் கோடலின் மனநலம் மோசமடைந்தது. 1977-இல் அவரது மனைவி
மருத்துவமனையில் சேர்க்கப்பட்டதால், அவருக்காக உணவு சமைக்க முடியாமல்
போனது. வாழ்நாள் முழுவதும் மனநலப் பிரச்சினைகளால் பாதிக்கப்பட்டிருந்த
கோடலுக்குத் தம்மைத் துன்புறுத்த யாரோ சதி செய்கிறார்கள் என்ற அச்சம்
மேலோங்கியது. நஞ்சூட்டப்படுவோம் என்று கடுமையாக அஞ்சியதால் உணவு உண்ண
மறுத்தார். 1978 ஜனவரி 14 அன்று பிரின்ஸ்டனில் பட்டினியால் இறந்தார்.

% SEGMENT OLP-0625-S06
\begin{reading}
கோடலின் முழுமையான வாழ்க்கை வரலாற்றுக்கு \citet{Dawson1997}-ஐப்
பார்க்கவும். மேலும் வாழ்க்கை வரலாற்றுக் கட்டுரைகளுக்கும் தருக்கம்,
மெய்யியல் ஆகியவற்றில் கோடலின் பங்களிப்புகள் பற்றிய கட்டுரைகளுக்கும்
\citet{Wang1990}, \citet{Baaz2011}, \citet{Takeuti2003},
\citet{Sigmund2007} ஆகியவற்றைப் பார்க்கவும்.

கோடலின் முனைவர் ஆய்வுரை மூல ஜெர்மன் மொழியில் கிடைக்கிறது
\citep{Godel1929}. முழுமையின்மைத் தேற்றங்களின் மூலக் கட்டுரை
\citep{Godel1931} ஆகும். வெளியிடப்பட்ட, வெளியிடப்படாத கோடலின்
எழுத்துகள் அனைத்தும், தேர்ந்தெடுக்கப்பட்ட கடிதங்களும், ஆங்கிலத்தில்
\emph{Collected Papers} என்னும் தொகுப்பில் கிடைக்கின்றன
\citet{Godel1986,Godel1990}.

கோடலின் முழுமையின்மைத் தேற்றங்கள் பற்றிய விரிவான விளக்கத்துக்கு
\citet{Smith2013}-ஐப் பார்க்கவும். தேற்றங்களைப் பற்றிய முறைசாராத,
மெய்யியல் நோக்கிலான உரையாடலுக்கு மார்க் லின்சன்மேயரின் பாட்காஸ்ட்டைப்
பார்க்கவும் \citep{Linsenmayer2014}.
\end{reading}

\end{document}
